\chapter{Materiales compuestos, laminados y uniones}\label{ch:09}
\objetivos{Relacionar orientación y transferencia de carga con las propiedades; utilizar reglas de mezcla como modelos limitados e identificar modos de fallo diferenciados.}
\section{Matriz, refuerzo e interfaz}
Un compuesto reúne fases distinguibles para obtener una respuesta que no se describe solo con la suma de sus masas. En un laminado de fibras largas, el refuerzo puede aportar gran rigidez en su dirección, mientras la matriz transmite cortadura, mantiene posición, protege y participa en cargas transversales. Decir que «la matriz nunca soporta carga» es incorrecto. Su contribución cambia con dirección, temperatura, geometría y estado de daño.

Las matrices poliméricas pueden ser termoestables o termoplásticas; su respuesta depende de temperatura, humedad y tiempo. Las matrices metálicas y cerámicas permiten otros intervalos térmicos y mecanismos de transferencia y daño. Un compuesto cerámico con mecanismos de desviación de grieta y extracción de fibras puede tener una respuesta al daño distinta de una cerámica monolítica.
\section{Anisotropía y modelos de mezcla}
Para un compuesto unidireccional ideal con fibras continuas alineadas, unión perfecta y deformación longitudinal igual en ambas fases,
\begin{equation}E_L=V_f E_f+V_m E_m,\qquad V_f+V_m=1.\label{eq:mezclaL}\end{equation}
Se obtiene sumando las fuerzas de fibras y matriz a igual deformación. Para una idealización de igual tensión en serie,
\begin{equation}\frac1{E_T}=\frac{V_f}{E_f}+\frac{V_m}{E_m}.\label{eq:mezclaT}\end{equation}
La segunda expresión es una idealización sencilla y no una predicción exacta de cualquier módulo transversal real. La geometría, las interfaces, vacíos y distribución del refuerzo hacen necesarios modelos o ensayos más representativos. La densidad ideal sin vacíos sí se calcula por promedio volumétrico $\rho=V_f\rho_f+V_m\rho_m$.
\caso{M09. Dirección de las fibras}{Para datos ficticios $V_f=0.60$, $E_f=230$ y $E_m=3\,\mathrm{GPa}$, el modelo longitudinal da $E_L=139.2\,\mathrm{GPa}$. La idealización en serie da $E_T=7.356\,\mathrm{GPa}$. La gran diferencia ilustra anisotropía, no dos resultados rivales para la misma dirección. El resultado transversal debe interpretarse dentro de la idealización, sin atribuirlo a una lámina comercial.
\par\textbf{Desarrollo para comparar tu intento.} La fracción de matriz es $V_m=1-V_f$. En paralelo longitudinal las fases comparten deformación y se suman rigideces ponderadas; en serie transversal comparten tensión y se suman flexibilidades antes de invertir.

\begin{enumerate}

\item Obtener fracción de matriz.

\[\fitmath{\mathrm{Vm}=1-0{,}6\approx 0{,}4\,}\]

\item Sumar las contribuciones longitudinales en las mismas unidades.

\[\fitmath{\mathrm{EL}=0{,}6\cdot 230+0{,}4\cdot 3\approx 139{,}2\,\mathrm{GPa}}\]

\item Sumar flexibilidades transversales y tomar el inverso.

\[\fitmath{\mathrm{ET}=\frac{1}{\frac{0{,}6}{230}+\frac{0{,}4}{3}}\approx 7{,}356077\,\mathrm{GPa}}\]

\end{enumerate}

}
\section{De la lámina al laminado}
Una lámina ortótropa necesita varios parámetros independientes: módulos, coeficientes de Poisson y módulos de cortadura, con relaciones de reciprocidad en elasticidad. Al girar la lámina, cambia la relación entre las componentes de tensión y deformación en los ejes globales. En un apilado, orientación y posición respecto al plano medio determinan respuesta extensional, de flexión y acoplamientos.

La teoría clásica de laminados organiza resultantes como
\begin{equation}\begin{bmatrix}\mathbf N\\\mathbf M\end{bmatrix}=\begin{bmatrix}\mathbf A&\mathbf B\\\mathbf B&\mathbf D\end{bmatrix}\begin{bmatrix}\boldsymbol\varepsilon^0\\\boldsymbol\kappa\end{bmatrix}.\label{eq:abd}\end{equation}
Un laminado simétrico respecto al plano medio tiene $\mathbf B=0$ en esta teoría y sin otros efectos que rompan las hipótesis. «Simétrico» no significa necesariamente «isótropo». Un apilado equilibrado elimina determinados acoplamientos de cortadura, pero tampoco constituye una solución universal frente a todas las cargas.
\ilustra{laminados}{Orientaciones de lámina, secuencia simétrica y estructura sándwich. La representación es conceptual, sin espesores de un producto real.}
\section{Fallo, delaminación y daño por impacto}
Los mecanismos incluyen rotura de fibras, fisuración de matriz, despegue de interfaz, inestabilidad local y delaminación. La resistencia puede depender de tracción o compresión y de dirección. Un criterio de fallo debe declarar qué mecanismos representa. Un único «límite elástico equivalente» de un metal no describe por sí solo una lámina anisótropa.

La delaminación separa capas y puede reducir mucho la capacidad de compresión o flexión sin una abertura superficial evidente. Las tensiones interlaminares se concentran en bordes, uniones y cambios de espesor; un modelo de tensión plana de lámina no las resuelve completamente. El daño apenas visible por impacto exige una estrategia de detección y evaluación compatible.
\section{Sándwich, GLARE y materiales de motor}
Una estructura sándwich separa caras resistentes mediante un núcleo ligero. La separación aumenta la rigidez a flexión con poca masa, pero introduce cortadura del núcleo, aplastamiento, arrugamiento de caras y despegue como modos posibles. No se verifica solo con la tensión longitudinal de las caras.

GLARE es un laminado híbrido de capas metálicas de aluminio y láminas reforzadas con fibra de vidrio. Puede combinar mecanismos de puenteo de grieta y respuesta a fatiga con exigencias de uniones y fabricación. Sus propiedades dependen del grado y apilado; no se le atribuye automáticamente la rigidez de un compuesto de carbono.

En motores pueden emplearse polímeros reforzados en zonas compatibles con su temperatura y ambiente, y compuestos de matriz cerámica en determinadas piezas calientes. Los CMC necesitan considerar oxidación, vapor de agua, recubrimientos ambientales, interfaces y vida; no son «cerámicas que ya no fallan de forma frágil» en cualquier condición.
\section{Uniones adhesivas}
Se distingue \textbf{fallo adhesivo}, situado en la interfaz entre adhesivo y sustrato; \textbf{fallo cohesivo}, dentro del adhesivo; y \textbf{fallo del sustrato}, dentro del material unido. La preparación superficial, contaminación, curado y diseño del extremo influyen en la distribución de esfuerzos.

Para una unión solapada, $F/(bL)$ es una tensión cortante media, no el máximo local. Los extremos pueden presentar concentraciones y esfuerzos de pelado. Aumentar indefinidamente el solape no garantiza un aumento proporcional de resistencia. Las reparaciones adheridas requieren datos, preparación y controles específicos.
\section{Desmontar las trampas habituales de compuestos y adhesivos}
La matriz soporta parte de las cargas y transmite esfuerzos entre fibras: no funciona simplemente como separador pasivo. Termoplásticos y termoestables son familias distintas; un termoestable curado y reticulado no se funde y reprocesa como un termoplástico. La elección debe atender a temperatura de transición vítrea, humedad, módulo y duración de la carga.

GLARE es un laminado híbrido metal--fibra de vidrio, cuya rigidez depende de secuencia de apilado y dirección. No existe un único módulo de GLARE válido para cualquier orientación. En un panel sándwich en flexión, las caras soportan gran parte de la tensión normal de tracción/compresión y el núcleo contribuye especialmente a cortadura y estabilidad de las caras. Los signos de las tensiones cambian con el momento, no se asignan por definición a una cara universal.

En uniones adheridas, fallo \emph{adhesivo} es separación interfacial; fallo \emph{cohesivo}, rotura dentro del adhesivo; y fallo \emph{del sustrato}, rotura del material unido. La resistencia real depende de preparación superficial, espesor, curado, dirección de carga y condiciones ambientales.


\section{De una lámina a un laminado: qué se suma}
Una lámina ortótropa tiene rigideces distintas según dirección y acoplamiento por Poisson. Al girar la lámina se transforma su matriz de rigidez; no se promedian simplemente módulos longitudinal y transversal. En teoría clásica de laminados, bajo sus hipótesis de placa delgada y unión entre capas, resultantes y deformaciones se relacionan mediante
\begin{equation}\begin{bmatrix}\mathbf N\\\mathbf M\end{bmatrix}=\begin{bmatrix}\mathbf A&\mathbf B\\\mathbf B&\mathbf D\end{bmatrix}\begin{bmatrix}\boldsymbol\varepsilon^0\\\boldsymbol\kappa\end{bmatrix}.\end{equation}
$\mathbf A$ reúne rigidez de membrana, $\mathbf D$ la de flexión y $\mathbf B$ el acoplamiento entre ambas. Las unidades son, respectivamente, fuerza por longitud, fuerza por longitud multiplicada por longitud al cuadrado, y fuerza, si se usan las resultantes habituales. La convención de signos de curvatura y del eje de espesor se declara antes de comparar matrices.

Un apilado simétrico respecto al plano medio anula $\mathbf B$ en el modelo ideal porque las contribuciones impares en espesor se cancelan. Eso no elimina todos los acoplamientos posibles dentro de $\mathbf A$ o $\mathbf D$ ni demuestra isotropía. Cambiar el orden de capas puede conservar una rigidez de membrana y modificar flexión. Las propiedades y espesores reales siguen siendo necesarios.
\section{Cargas de matriz, interfaces y daño}
Las fibras pueden dominar una rigidez longitudinal, pero la matriz transmite carga entre fibras, estabiliza y participa en respuesta transversal y cortante. La interfaz determina transferencia y ciertos modos de daño. Bajo compresión, flexión o carga fuera de plano, afirmar que la matriz «no soporta carga» puede ocultar el mecanismo que gobierna.

Fisuración de matriz, rotura de fibra y delaminación pueden aparecer en secuencia e interactuar. Un impacto con huella pequeña puede dejar daño interno relevante; observar la superficie no caracteriza todo el volumen. La resistencia interlaminar y la tenacidad de delaminación describen aspectos distintos y dependen del modo de apertura o cortante. La teoría clásica elástica de placa no resuelve por sí sola el campo tridimensional de un borde libre o un impacto.
\Needspace{.48\textheight}
\section{Sándwich, GLARE y uniones adhesivas}
Separar caras aumenta rigidez a flexión con poca cantidad de material resistente de cara, pero el núcleo debe transmitir cortante y mantener la separación. Arrugamiento, aplastamiento, despegue e infiltración pueden gobernar antes de alcanzar la resistencia nominal de las caras. La ganancia ideal de inercia del problema P25 debe leerse junto a estas limitaciones.

GLARE combina capas metálicas y de compuesto de vidrio en una arquitectura específica; no es cualquier panel con vidrio y aluminio. Las rutas de daño, propagación y reparación dependen de esa arquitectura. En una unión adhesiva, $F/(bL)$ solo es cortante media: extremos del solape y excentricidad introducen concentraciones y pelado. Preparación superficial, curado, espesor, ambiente y compatibilidad de sustratos son parte del sistema resistente.
\repaso{¿Un laminado simétrico y equilibrado es necesariamente isótropo?}{No. Simetría y equilibrio eliminan ciertos acoplamientos bajo condiciones del modelo, pero las rigideces pueden seguir variando con dirección. La isotropía o cuasiisotropía requiere condiciones adicionales sobre el apilado y el comportamiento considerado.}

\resumen{La dirección de carga forma parte de la propiedad. Las reglas de mezcla son modelos con hipótesis, no fichas de laminado. Simetría del apilado y equilibrio no equivalen a isotropía. Adhesivo, cohesivo y sustrato identifican localizaciones distintas del fallo.}
\fuentes{Roylance, \textit{Introduction to Composites} y \textit{Laminated Composite Plates} \cite{MIT-COMP,MIT-LAM}; FAA--General, cap. 7 \cite{FAA-G}; NASA sobre CMC \cite{NASA-CMC}.}
